% TOP_PROBLEM: {"id": 1364, "title": "Cap Set Problem", "category_id": 2, "category": {"id": 2, "name": "combinatorics", "display_name": "Combinatorics", "description": "Counting problems, graph theory, discrete structures.", "slug": "combinatorics", "order_index": 2, "created_at": "2026-07-31T15:26:25.671Z"}, "status": "open"}
% FIRST_POSTED: null
% ENTRY_KIND: statement_only
% Imported from: batch06.tex; UnsolvedMath ID 1364
\documentclass[11pt]{article}
\usepackage[margin=25mm]{geometry}
\usepackage{fontspec}
\setmainfont{Latin Modern Roman}
\usepackage{amsmath,amssymb,mathtools}
\usepackage{unicode-math}
\setmathfont{Latin Modern Math}
\usepackage{xurl,hyperref}
\hypersetup{colorlinks=true,urlcolor=blue}
\setcounter{secnumdepth}{0}
\setlength{\parindent}{0pt}
\setlength{\parskip}{5pt}
\setlength{\emergencystretch}{3em}
\newcommand{\sref}[2]{\hyperlink{source:#1:#2}{[S#2]}}
\newcommand{\cataloguescope}[1]{\paragraph{Recorded target.}#1}
\begin{document}
% UnsolvedMath ID 1364; source catalogue code GAME-007
\setcounter{section}{1363}
\section{Cap Set Problem}\label{1364}
{\small\sffamily UnsolvedMath ID 1364 · Combinatorics}
\subsection{Definitions and mathematical statement}

\paragraph{Shared notation.}$\mathbb N=\{1,2,\ldots\}$, and $\mathbb Z,\mathbb Q,\mathbb R,\mathbb C$ denote the integers, rationals, reals and complex numbers. Any other conventions are specified locally.


Let $\mathbb F_3=\mathbb Z/3\mathbb Z$. A cap set is a subset $A\subseteq\mathbb F_3^n$ such that $x+y+z=0$ with $x,y,z\in A$ implies $x=y=z$. This is equivalent to containing no affine line $\{a,a+v,a+2v\}$ with $v\ne0$. Let
\[
 m(n)=\max\{|A|:\ A\subseteq\mathbb F_3^n\text{ is a cap set}\}.
\]
\paragraph{Statement recorded in the supplied source.}
Determine $m(n)$ for general $n\in\mathbb N$, and its asymptotically sharp growth as $n\to\infty$. In particular, determine its exponential growth rate $\limsup_{n\to\infty}m(n)^{1/n}$. \sref{1364}{1}\sref{1364}{2}
\subsection{Short English statement}
Find the largest subset of each n-dimensional affine space over the three-element field containing no three distinct collinear points, and determine its growth with dimension.
\subsection{Sources}
\begin{description}
\item[\hypertarget{source:1364:1}{S1}] ULAM.AI and the UnsolvedMath Contributors, \emph{UnsolvedMath: A Curated Collection of Open Mathematics Problems}, record 1364 (GAME-007). CC BY 4.0. \url{https://huggingface.co/datasets/ulamai/UnsolvedMath}.
\item[\hypertarget{source:1364:2}{S2}] Fred Tyrrell, \emph{New Lower Bounds for Cap Sets}, Discrete Analysis 2023:20; arXiv:2209.10045v2, Definition 1.1 and Section 1. \url{https://arxiv.org/abs/2209.10045v2}.
\end{description}
\label{end:1364}
\end{document}
